% SOURCE: Chinese PDF viewer pp. 104--109. Section 3.
\section{Degenerate Perturbation Theory}
Suppose
\begin{equation}
 \op H^{(0)}\ket{n,i}^{(0)}=E_n^{(0)}\ket{n,i}^{(0)},
 \qquad i=1,2,\ldots,a.
\end{equation}
The zeroth-order energy has $a$-fold degeneracy. Schmidt orthogonalization or the introduction of a complete set of mechanical quantities can guarantee mutual orthonormality:
\begin{equation}
 {}^{(0)}\!\braket{n,i}{n,j}^{(0)}=\delta_{ij}.
\end{equation}
\textbf{Question:} Which one of the states $\ket{n,i}^{(0)}$ should be chosen as the wavefunction for $E_n^{(0)}$?

Take the general form
\begin{equation}
 \ket n^{(0)}=\sum_{i=1}^{a}c_i\ket{n,i}^{(0)}.
\end{equation}
It is still an eigenstate of $\op H^{(0)}$ belonging to $E_n^{(0)}$.

\textbf{Question:} How should the $c_i$ be chosen?

Substitute $\ket n^{(0)}$ into the first-order approximation equation:
\begin{equation}
 (\op H^{(0)}-E_n^{(0)})\ket n^{(1)}
 =-(\op H^{(1)}-E_n^{(1)})\ket n^{(0)},
\end{equation}
so that
\begin{equation}
 (\op H^{(0)}-E_n^{(0)})\ket n^{(1)}
 =-\sum_{i=1}^a c_i(\op H^{(1)}-E_n^{(1)})\ket{n,i}^{(0)}.
\end{equation}
Multiply on the left by ${}^{(0)}\!\bra{n,j}$:
\begin{equation}
 0=\sum_{i=1}^{a}c_i(H_{ji}^{(1)}-E_n^{(1)}\delta_{ji}),
 \qquad H_{ji}^{(1)}={}^{(0)}\!\mel{n,j}{\op H^{(1)}}{n,i}^{(0)}.
\end{equation}
This is a homogeneous system of equations for the coefficients $c_i$ ($a$ equations, $j=1,\ldots,a$, determining $a$ coefficients, $i=1,\ldots,a$). For a nonzero solution, its determinant must vanish:
\begin{equation}
 \begin{vmatrix}
 H_{11}^{(1)}-E_n^{(1)}&H_{12}^{(1)}&\cdots&H_{1a}^{(1)}\\
 H_{21}^{(1)}&H_{22}^{(1)}-E_n^{(1)}&\cdots&H_{2a}^{(1)}\\
 \vdots&\vdots&\ddots&\vdots\\
 H_{a1}^{(1)}&H_{a2}^{(1)}&\cdots&H_{aa}^{(1)}-E_n^{(1)}
 \end{vmatrix}=0.
\end{equation}
The secular equation therefore determines $a$ roots $E_{ni}^{(1)}$ of the first-order energy correction:
\begin{equation}
 E_n^{(0)}\longrightarrow E_{ni}=E_n^{(0)}+E_{ni}^{(1)},
 \qquad i=1,\ldots,a.
\end{equation}
If the $a$ roots of the secular equation are distinct, the degeneracy is completely removed; if roots repeat, it is only partially removed.

\textbf{Note.} If $\op H^{(1)}$ is diagonal in the degenerate subspace, $E_{ni}^{(1)}=H_{ii}^{(1)}$.

Substituting each $E_{ni}^{(1)}$ back into the homogeneous equations gives a set of coefficients $\{c_{i1},c_{i2},\ldots,c_{ia}\}$, and hence the zeroth-order wavefunction of $\op H$ corresponding to $E_{ni}$:
\begin{equation}
 E_{ni}=E_n^{(0)}+E_{ni}^{(1)},\qquad
 \ket{n_i}^{(0)}=\sum_{j=1}^{a}c_{ij}\ket{n,j}^{(0)}.
\end{equation}
Higher-order corrections can be found in the references.

\textbf{Example 1: A doubly degenerate system.}
Let $E_n^{(0)}$ have basis states $\ket{n,1}^{(0)}$ and $\ket{n,2}^{(0)}$. The secular equation in this two-dimensional subspace is
\begin{equation}
 \begin{vmatrix}
 H_{11}^{(1)}-E_n^{(1)}&H_{12}^{(1)}\\
 H_{21}^{(1)}&H_{22}^{(1)}-E_n^{(1)}
 \end{vmatrix}=0,
 \quad
 H_{ij}^{(1)}={}^{(0)}\!\mel{n,i}{\op H^{(1)}}{n,j}^{(0)},\quad i,j=1,2.
\end{equation}
Thus
\begin{equation}
 (E_n^{(1)})^2-(H_{11}^{(1)}+H_{22}^{(1)})E_n^{(1)}
 +H_{11}^{(1)}H_{22}^{(1)}-|H_{12}^{(1)}|^2=0,
\end{equation}
\begin{equation}
 E_n^{(1)}=\frac12\left[
 H_{11}^{(1)}+H_{22}^{(1)}
 \pm\sqrt{(H_{11}^{(1)}-H_{22}^{(1)})^2+4|H_{12}^{(1)}|^2}
 \right].
\end{equation}
If $H_{12}^{(1)}=0$, then
\begin{equation}
 E_{n1}^{(1)}=H_{11}^{(1)},\quad E_{n2}^{(1)}=H_{22}^{(1)},\quad
 E_{n1}=E_n^{(0)}+H_{11}^{(1)},\quad E_{n2}=E_n^{(0)}+H_{22}^{(1)}.
\end{equation}

\textbf{Example 2: The hydrogen atom in an external electric field (Stark effect).}
Choose the $z$ axis along the external electric field $\vect{\mathcal E}$:
\begin{equation}
 \op H=-\frac{\hbar^2}{2\mu}\nabla^2-\frac{e^2}{r}
 +e\vect{\mathcal E}\cdot\vect r
 =-\frac{\hbar^2}{2\mu}\nabla^2-\frac{e^2}{r}
 +e\mathcal E r\cos\theta.
\end{equation}
For a weak field $\mathcal E$,
\begin{equation}
 \op H^{(0)}=-\frac{\hbar^2}{2\mu}\nabla^2-\frac{e^2}{r},\qquad
 \op H^{(1)}=e\mathcal E r\cos\theta.
\end{equation}
Zeroth-order approximation:
\begin{equation}
 E_n^{(0)}=-\frac{\mu e^4}{2n^2\hbar^2},\qquad
 \psi_{nlm}(\vect r),\quad l=0,\ldots,n-1,\quad m=-l,\ldots,l,
\end{equation}
with degeneracy
\begin{equation}
 g=\sum_{l=0}^{n-1}(2l+1)=n^2.
\end{equation}
For $n=2$, the zeroth-order energy $E_2^{(0)}$ is fourfold degenerate, with states
\begin{equation}
 \phi_1=\psi_{200},\qquad \phi_2=\psi_{210},\qquad
 \phi_3=\psi_{211},\qquad \phi_4=\psi_{21,-1}.
\end{equation}
To find the first-order splitting, calculate the matrix elements
\begin{equation}
 H_{ij}^{(1)}=\int\dd^3r\,\phi_i^*(\vect r)\,
 \op H^{(1)}\phi_j(\vect r),\qquad i,j=1,2,3,4.
\end{equation}
The matrix element is
\begin{equation}
 H_{12}^{(1)}=H_{21}^{(1)}=-3e\mathcal E a_0,
\end{equation}
with all other entries zero. The secular equation in the four-dimensional subspace is
\begin{equation}
 \begin{vmatrix}
 -E_2^{(1)}&-3e\mathcal E a_0&0&0\\
 -3e\mathcal E a_0&-E_2^{(1)}&0&0\\
 0&0&-E_2^{(1)}&0\\
 0&0&0&-E_2^{(1)}
 \end{vmatrix}=0.
\end{equation}
Hence
\begin{equation}
 E_{21}^{(1)}=3e\mathcal E a_0,\qquad
 E_{22}^{(1)}=-3e\mathcal E a_0,\qquad
 E_{23}^{(1)}=E_{24}^{(1)}=0.
\end{equation}
The original fourfold-degenerate level splits into one level shifted upward by $3e\mathcal E a_0$, a twofold-degenerate level with zero shift, and one level shifted downward by $3e\mathcal E a_0$. The degenerate perturbation is partially removed.

\textbf{Example 3: Relativistic correction for the hydrogen atom.}
The nonrelativistic kinetic energy is
\begin{equation}
 T=\frac{p^2}{2m}.
\end{equation}
The relativistic kinetic energy is
\begin{equation}\begin{aligned}
 T&=\sqrt{p^2c^2+m^2c^4}-mc^2
 =mc^2\sqrt{1+\frac{p^2}{m^2c^2}}-mc^2\\
 &=\frac{p^2}{2m}-\frac{p^4}{8m^3c^2}+\cdots.
\end{aligned}\end{equation}
Keeping only the $p^4$ correction gives the relativistic Hamiltonian
\begin{equation}
 \op H=\sqrt{\op p^{\,2}c^2+m^2c^4}-mc^2-\frac{e^2}{r}
 \simeq \frac{\op p^{\,2}}{2m}-\frac{e^2}{r}
 -\frac{\op p^{\,4}}{8m^3c^2}.
\end{equation}
Take
\begin{equation}
 \op H^{(0)}=\frac{\op p^{\,2}}{2m}-\frac{e^2}{r},\qquad
 \op H^{(1)}=-\frac{\op p^{\,4}}{8m^3c^2}.
\end{equation}
The zeroth-order energy is $E_n^{(0)}$ and the $n^2$ eigenstates are $\ket{nlm}$, with wavefunctions
\begin{equation}
 \braket{\vect r}{nlm}=\psi_{nlm}(\vect r)=R_{nl}(r)Y_{lm}(\theta,\varphi).
\end{equation}
In the $n^2$-dimensional subspace,
\begin{equation}
 H_{nlm,n'l'm'}^{(1)}=-\frac{1}{8m^3c^2}
 \mel{nlm}{\op p^{\,4}}{n'l'm'}.
\end{equation}
The zeroth-order Schr\"odinger equation is
\begin{equation}
 \left(\frac{\op p^{\,2}}{2m}+V(r)\right)\ket{nlm}
 =E_n^{(0)}\ket{nlm},
 \qquad
 \op p^{\,2}\ket{nlm}=2m(E_n^{(0)}-V(r))\ket{nlm}.
\end{equation}
Using this relation, the matrix elements become
\begin{equation}\begin{aligned}
 H_{nlm,n'l'm'}^{(1)}
 &=-\frac{1}{2mc^2}\mel{nlm}{(E_n^{(0)}-V(r))^2}{n'l'm'}\\
 &=-\frac{1}{2mc^2}\left[
 (E_n^{(0)})^2\delta_{nn'}\delta_{ll'}\delta_{mm'}
 -2E_n^{(0)}\mel{nlm}{V(r)}{n'l'm'}
 +\mel{nlm}{V^2(r)}{n'l'm'}\right].
\end{aligned}\end{equation}
Since $V(r)$ depends only on $r$ and not on the angles,
\begin{equation}
 \mel{nlm}{V(r)}{n'l'm'}
 =\int\dd^3r\,R_{nl}(r)Y_{lm}^*(\theta,\varphi)
 \left(-\frac{e^2}{r}\right)R_{n'l'}(r)Y_{l'm'}(\theta,\varphi)
 =-\frac{e^2}{n^2a_0}\delta_{ll'}\delta_{mm'}
\end{equation}
within a fixed principal-energy subspace. Similarly,
\begin{equation}
 \mel{nlm}{V^2(r)}{nl'm'}
 =\frac{e^4}{(l+1/2)n^3a_0^2}\delta_{ll'}\delta_{mm'}.
\end{equation}
Therefore the first-order matrix is diagonal:
\begin{equation}
 H_{nlm,nl'm'}^{(1)}=H_{nl}^{(1)}\delta_{ll'}\delta_{mm'},
\end{equation}
\begin{equation}
 H_{nl}^{(1)}=-\frac{1}{2mc^2}
 \left[(E_n^{(0)})^2+\frac{2e^2}{n^2a_0}E_n^{(0)}
 +\frac{e^4}{(l+1/2)n^3a_0^2}\right].
\end{equation}
This shows that the eigenstates of $\op H^{(0)}$ are also eigenstates of $\op H^{(1)}$ in this subspace, and the secular equation reduces to $E_{nl}^{(1)}=H_{nl}^{(1)}$:
\begin{equation}
 E_n^{(0)}\quad(n^2\text{-fold degenerate})
 \longrightarrow E_{nl}=E_n^{(0)}+H_{nl}^{(1)}
 \quad((2l+1)\text{-fold degenerate}).
\end{equation}

\textbf{Example 4: The fine structure of the spectrum.}
Consider the spin--orbit coupling of the electron:
\begin{equation}
 \op H=\op H^{(0)}+\op H^{(1)},\qquad
 \op H^{(0)}=-\frac{\hbar^2}{2\mu}\nabla^2+V(r),
\end{equation}
\begin{equation}
 \op H^{(1)}=\frac{1}{2\mu^2r}\frac{\dd V(r)}{\dd r}
 \op{\vect L}\cdot\op{\vect S}
 \equiv\xi(r)\op{\vect L}\cdot\op{\vect S}.
\end{equation}
Zeroth-order states: the system is in an energy eigenstate $E_{nl}^{(0)}$. For a given $l$, the energy level has $2(2l+1)$ degeneracy. For this level the basis may be chosen as
\begin{equation}
 \{\ket{n,l,m,s_z}\},\qquad
 \ket{n,l,j,m_j},
\end{equation}
using either the $\{\op H^{(0)},\op{\vect L}^{\,2},\op L_z,
 \op{\vect S}^{\,2},\op S_z\}$ representation or the
$\{\op H^{(0)},\op{\vect J}^{\,2},\op{\vect L}^{\,2},\op J_z\}$ representation. The latter describes an eigenstate $\ket{nljm_j}$.

The zeroth-order energy is degenerate, so find the matrix elements of $\op H^{(1)}$. Which representation should be used?

In the uncoupled representation, $[\op H^{(1)},\op L_z]\ne0$ and $[\op H^{(1)},\op S_z]\ne0$, so $\ket{nlms_z}$ is not an eigenstate of $\op H^{(1)}$. Its matrix elements need not be diagonal and the secular equation is more complicated.

In the coupled representation, $[\op H^{(1)},\op{\vect J}^{\,2}]=0$ and $[\op H^{(1)},\op J_z]=0$. Thus $\ket{nljm_j}$ is an eigenstate of $\op H^{(1)}$, whose matrix is diagonal. The first-order energy correction is
\begin{equation}\begin{aligned}
 E_{nlj}^{(1)}&=\mel{nljm_j}{\op H^{(1)}}{nljm_j}\\
 &=\mel{nljm_j}{\xi(r)\op{\vect L}\cdot\op{\vect S}}{nljm_j}\\
 &=\mel{nljm_j}{\xi(r)\frac12
 (\op{\vect J}^{\,2}-\op{\vect L}^{\,2}-\op{\vect S}^{\,2})}{nljm_j}\\
 &=\frac{\hbar^2}{2}[j(j+1)-l(l+1)-s(s+1)]
 \int_0^\infty R_{nl}^{\,2}(r)\xi(r)r^2\dd r.
\end{aligned}\end{equation}
The level changes as
\begin{equation}
 E_{nl}^{(0)}\quad(2(2l+1)\text{-fold degenerate})
 \longrightarrow E_{nlj}=E_{nl}^{(0)}+E_{nlj}^{(1)}
 \quad((2j+1)\text{-fold degenerate}).
\end{equation}
For example, when $n=2$,
\begin{equation}
 E_2^{(0)}\quad(8\text{-fold degenerate})\ \longrightarrow\
 \begin{cases}
 E_2^{(0)}+E_{2,1,3/2}^{(1)} & (4\text{-fold}),\\
 E_2^{(0)}+E_{2,1,1/2}^{(1)} & (2\text{-fold}),\\
 E_2^{(0)}+E_{2,0,1/2}^{(1)} & (2\text{-fold}).
 \end{cases}
\end{equation}
\textbf{Note:} when $l=0$, $j=s$ and $E_{2,0,1/2}^{(1)}=0$.